\chapter{Complexity — P Subset P Poly Gadget}
%%% generated by the blueprint pipeline from TCSlib.Complexity.CircuitComplexity.PSubsetPPolyGadget
\section{Overview}

Declarations of the Lean module \texttt{TCSlib.Complexity.CircuitComplexity.PSubsetPPolyGadget}.

\section{Declarations}

\begin{definition}[Boolean function circuit gadget]
\lean{BoolCircuit.gadget}
\label{BoolCircuit.gadget}
\leanok
\uses{BoolCircuit.DAGCircuit, BoolCircuit.DAGCircuit.binarize, BoolCircuit.TreeCircuit.toDAG, BoolCircuit.TreeCircuit.toDAG_isWellFormed, BoolCircuit.universalCircuit}
For a nonnegative integer $m$ and a Boolean function $f:\{0,1\}^m\to\{0,1\}$, the gadget
is the directed acyclic circuit on $m$ input vertices obtained by converting the
universal circuit for $f$ to a directed acyclic graph and binarizing it.
\end{definition}

\begin{theorem}[Correct evaluation of a Boolean-function gadget]
\lean{BoolCircuit.gadget_eval}
\label{BoolCircuit.gadget_eval}
\leanok
\uses{BoolCircuit.DAGCircuit.binarize_eval, BoolCircuit.DAGCircuit.eval, BoolCircuit.TreeCircuit.toDAG_eval, BoolCircuit.gadget, BoolCircuit.universalCircuit_eval}
\difficulty{3}
For every natural number $m$, every Boolean function $f$ on assignments to $m$
variables, and every assignment $x$ to those variables, the output of the gadget
constructed from $f$ on input $x$ equals $f(x)$.
\end{theorem}

\begin{theorem}[Gadget has fan-in at most two]
\lean{BoolCircuit.gadget_isFaninTwo}
\label{BoolCircuit.gadget_isFaninTwo}
\leanok
\uses{BoolCircuit.DAGCircuit.IsFaninTwo, BoolCircuit.DAGCircuit.binarize_isFaninTwo, BoolCircuit.gadget}
\difficulty{1}
For every natural number $m$ and every Boolean function $f$ on assignments to $m$
variables, the gadget constructed from $f$ is well formed and every gate in it has at
most two inputs.
\end{theorem}

\begin{definition}[Deduplicated gate arguments]
\lean{BoolCircuit.DAGGate.dedupArgs}
\label{BoolCircuit.DAGGate.dedupArgs}
\leanok
\uses{BoolCircuit.DAGGate}
For a gate with a specified kind and an ordered list of argument vertices, its
deduplicated version has the same kind and the list obtained by removing repeated
arguments.
\end{definition}

\begin{theorem}[Evaluation after removing repeated inputs]
\lean{BoolCircuit.DAGGate.eval_dedupArgs}
\label{BoolCircuit.DAGGate.eval_dedupArgs}
\leanok
\uses{BoolCircuit.DAGGate, BoolCircuit.DAGGate.dedupArgs, BoolCircuit.DAGGate.eval}
\difficulty{5}
Let a Boolean gate be an AND gate, an OR gate, or a NOT gate, with a finite list of
natural-number input indices, and let the available values be a finite list of Booleans
(an out-of-range index is read as false). Removing repeated indices from the gate’s
input list leaves its evaluation unchanged.
\end{theorem}

\begin{theorem}[Duplicate arguments do not affect evaluation]
\lean{BoolCircuit.runWith_eval_map_dedupArgs}
\label{BoolCircuit.runWith_eval_map_dedupArgs}
\leanok
\uses{BoolCircuit.DAGGate, BoolCircuit.DAGGate.dedupArgs, BoolCircuit.DAGGate.eval, BoolCircuit.DAGGate.eval_dedupArgs, BoolCircuit.runWith, BoolCircuit.runWith_cons}
\difficulty{4}
For any finite list of Boolean gates and any initial Boolean value list, successively
evaluating the gates after removing repeated entries from each gate's argument list
gives the same final list as evaluating the original gates.
\end{theorem}

\begin{theorem}[Mapped list reconstruction from indexed values]
\lean{BoolCircuit.ofFn_getD_eq_map}
\label{BoolCircuit.ofFn_getD_eq_map}
\leanok
\difficulty{4}
Let $l$ be a list of natural numbers of length $m$, and let $f$ map natural numbers to
Booleans. The list obtained by applying $f$ to the entries of $l$ is the list of values
$f(l_i)$ for indices $i=0,\ldots,m-1$, where each entry is read with default value $0$.
\end{theorem}

\begin{theorem}[Mapping Commutes with List Lookup]
\lean{BoolCircuit.getD_map_getD}
\label{BoolCircuit.getD_map_getD}
\leanok
\difficulty{3}
Let $l$ be a list of natural numbers, let $f$ map natural numbers to Booleans, and let
$i$ be an index less than the length of $l$. Applying $f$ to the entry of $l$ at index
$i$ gives the same Boolean as looking up index $i$ in the list obtained by applying $f$
to every entry of $l$, with default value false.
\end{theorem}

\begin{definition}[Vertex embedding map]
\lean{BoolCircuit.embedMap}
\label{BoolCircuit.embedMap}
\leanok
Given natural numbers $m,L$, a list $s$ of natural numbers, and a natural number $v$,
map $v$ to the entry $s_v$ when $v<m$, using $0$ if that entry is unavailable; when
$v\ge m$, map it to $L+(v-m)$.
\end{definition}

\begin{definition}[Embedded circuit gates]
\lean{BoolCircuit.embedGates}
\label{BoolCircuit.embedGates}
\leanok
\uses{BoolCircuit.DAGCircuit, BoolCircuit.DAGGate, BoolCircuit.DAGGate.dedupArgs, BoolCircuit.DAGGate.remap, BoolCircuit.embedMap}
Let $D$ be a directed acyclic circuit with $m$ input vertices, let $s$ be a list of
source vertices, and let $L$ be a natural number. Replace each gate argument of $D$ by
its image under the vertex map that sends $v<m$ to $s_v$ (or $0$ if unavailable) and
sends $v\ge m$ to $L+(v-m)$, removing repeated arguments while preserving their order.
Append one AND gate whose sole argument is the image of $D$'s output vertex under that
map.
\end{definition}

\begin{theorem}[Length of the embedded gate list]
\lean{BoolCircuit.length_embedGates}
\label{BoolCircuit.length_embedGates}
\leanok
\uses{BoolCircuit.DAGCircuit, BoolCircuit.embedGates}
\difficulty{1}
Let $D$ be a directed acyclic circuit with $m$ input vertices, let $srcs$ be a list of
natural numbers, and let $L$ be a natural number. The gate list formed by embedding $D$
using $srcs$ and $L$ has length equal to the number of gates of $D$ plus one.
\end{theorem}

\begin{theorem}[Bounded Source Index Lookup]
\lean{BoolCircuit.getD_src_lt}
\label{BoolCircuit.getD_src_lt}
\leanok
\difficulty{3}
Let $m$ and $L$ be natural numbers, and let $\mathrm{srcs}$ be a list of length $m$
whose entries are all less than $L$. For every $v<m$, the entry of $\mathrm{srcs}$ at
index $v$, using $0$ as the default out of range, is less than $L$.
\end{theorem}

\begin{theorem}[Bound for embedded vertices]
\lean{BoolCircuit.embedMap_lt}
\label{BoolCircuit.embedMap_lt}
\leanok
\uses{BoolCircuit.embedMap, BoolCircuit.getD_src_lt}
\difficulty{4}
Let $m,L$ be natural numbers, and let $\mathrm{srcs}$ be a list of length $m$ whose
entries are all less than $L$. For any natural numbers $v,i$ with $v<m+i$, the embedded
vertex obtained by sending vertices below $m$ to the corresponding entry of
$\mathrm{srcs}$ and vertices at least $m$ to $L+(v-m)$ is less than $L+i$.
\end{theorem}

\begin{theorem}[Evaluation of an embedded circuit]
\lean{BoolCircuit.runWith_embedGates_getD}
\label{BoolCircuit.runWith_embedGates_getD}
\leanok
\uses{BoolCircuit.DAGCircuit, BoolCircuit.DAGCircuit.eval, BoolCircuit.DAGGate.dedupArgs, BoolCircuit.DAGGate.eval, BoolCircuit.DAGGate.eval_remap, BoolCircuit.DAGGate.remap, BoolCircuit.embedGates, BoolCircuit.embedMap, BoolCircuit.embedMap_lt, BoolCircuit.runWith, BoolCircuit.runWith_append, BoolCircuit.runWith_eval_map_dedupArgs, BoolCircuit.runWith_remap_rel, BoolCircuit.runWith_singleton}
\difficulty{6}
Let $D$ be a directed acyclic Boolean circuit with $m$ inputs, let $s$ be a list of $m$
source indices, and let $I$ be a list of Boolean values such that every index in $s$ is
less than the length of $I$. Running the embedded gates of $D$ from $I$ yields at
position $|I|+|D.	ext{gates}|$ the value of $D$ on the input whose $i$th coordinate is
the value of $I$ at the index specified by the $i$th entry of $s$; an out-of-range
lookup in this input is interpreted as false.
\end{theorem}

\begin{theorem}[Acyclicity of an embedded circuit]
\lean{BoolCircuit.gatesAcyclic_embedGates}
\label{BoolCircuit.gatesAcyclic_embedGates}
\leanok
\uses{BoolCircuit.DAGCircuit, BoolCircuit.DAGGate.dedupArgs, BoolCircuit.DAGGate.remap, BoolCircuit.GatesAcyclic, BoolCircuit.GatesAcyclic.snoc, BoolCircuit.embedGates, BoolCircuit.embedMap, BoolCircuit.embedMap_lt}
\difficulty{5}
Let $D$ be a finite directed acyclic circuit with $m$ input vertices, and let
$\mathrm{srcs}$ be a list of $m$ vertex indices, each less than $L$. Embed the gates of
$D$ using these sources and offset $L$, and append the gate that takes the embedded
output as its sole input. The resulting gate list is acyclic relative to $L$.
\end{theorem}

\begin{theorem}[Fan-in of an embedded circuit]
\lean{BoolCircuit.faninTwo_embedGates}
\label{BoolCircuit.faninTwo_embedGates}
\leanok
\uses{BoolCircuit.DAGCircuit, BoolCircuit.DAGCircuit.IsFaninTwo, BoolCircuit.DAGGate.dedupArgs, BoolCircuit.DAGGate.remap, BoolCircuit.embedGates, BoolCircuit.embedMap}
\difficulty{5}
Let $m$ be a natural number, and let $D$ be a circuit with $m$ inputs that is well
formed and whose gates each have at most two inputs. For any list of source vertices and
any natural number $L$, every gate in the embedding of $D$ is well formed and has at
most two inputs.
\end{theorem}

\begin{definition}[Concatenated embeddings of circuits]
\lean{BoolCircuit.embedAll}
\label{BoolCircuit.embedAll}
\leanok
\uses{BoolCircuit.DAGCircuit, BoolCircuit.DAGGate, BoolCircuit.embedGates}
Let $m$ and $L$ be natural numbers, let $\mathcal D$ be a list of directed acyclic
circuits with $m$ input vertices, and let $s$ be a list of natural-number vertices. For
each circuit in order, append its embedded gate list using the same source list $s$,
starting with offset $L$ and increasing the offset after each circuit by its number of
gates plus one; the resulting list is the concatenation of these embedded gate lists.
\end{definition}

\begin{definition}[Total embedded width]
\lean{BoolCircuit.embedWidth}
\label{BoolCircuit.embedWidth}
\leanok
\uses{BoolCircuit.DAGCircuit}
For a list $\mathcal D$ of directed acyclic circuits, define its embedded width
recursively as the sum, over each circuit $D$ in the list, of its number of gates plus
one.
\end{definition}

\begin{definition}[Circuit embedding offset]
\lean{BoolCircuit.embedOffset}
\label{BoolCircuit.embedOffset}
\leanok
\uses{BoolCircuit.DAGCircuit}
For a list $\mathcal D$ of directed acyclic circuits and a natural number $j$, define
the offset recursively: it is $0$ for the empty list; for a nonempty list, it is the
first circuit's number of gates when $j=0$, and otherwise it is that number plus one
plus the offset at index $j-1$ in the remaining list.
\end{definition}

\begin{theorem}[Length of the combined embedded gates]
\lean{BoolCircuit.length_embedAll}
\label{BoolCircuit.length_embedAll}
\leanok
\uses{BoolCircuit.DAGCircuit, BoolCircuit.embedAll, BoolCircuit.embedWidth}
\difficulty{4}
For any finite list $D_1,\ldots,D_r$ of directed acyclic circuits with $m$ input
vertices, any list $srcs$ of natural numbers, and any natural number $L$, the number of
gates in the list formed by embedding these circuits successively equals the sum of
their widths, where the width of each circuit is its number of gates plus one.
\end{theorem}

\begin{theorem}[Embedding offset bound]
\lean{BoolCircuit.embedOffset_lt}
\label{BoolCircuit.embedOffset_lt}
\leanok
\uses{BoolCircuit.DAGCircuit, BoolCircuit.embedOffset, BoolCircuit.embedWidth}
\difficulty{4}
Let $m$ be a natural number, and let $D_1,\ldots,D_r$ be a list of circuits with $m$
inputs. For every index $j<r$, the offset assigned to the $j$th circuit is less than the
total embedding width of the list.
\end{theorem}

\begin{theorem}[Evaluation of an embedded circuit]
\lean{BoolCircuit.runWith_embedAll_getD}
\label{BoolCircuit.runWith_embedAll_getD}
\leanok
\uses{BoolCircuit.DAGCircuit, BoolCircuit.DAGCircuit.eval, BoolCircuit.DAGGate.eval, BoolCircuit.embedAll, BoolCircuit.embedGates, BoolCircuit.embedOffset, BoolCircuit.getD_src_lt, BoolCircuit.runWith, BoolCircuit.runWith_append, BoolCircuit.runWith_embedGates_getD, BoolCircuit.runWith_getD_of_lt}
\difficulty{6}
Let $D_1,\ldots,D_k$ be a finite list of directed acyclic Boolean circuits, each with $m$
inputs, and let $s$ be a list of $m$ indices into an initial Boolean list $I$. Assume
every index in $s$ is less than the length of $I$. For any $j<k$, after executing the
gates obtained by embedding all the circuits using $s$, starting from $I$, the value at
the index given by the length of $I$ plus the embedding offset of $D_{j+1}$ equals the
output of $D_{j+1}$ on the input whose $i$th coordinate is the value of $I$ at the index
specified by the $i$th entry of $s$, with out-of-range lookups interpreted as false.
\end{theorem}

\begin{theorem}[Acyclicity of an embedded circuit list]
\lean{BoolCircuit.gatesAcyclic_embedAll}
\label{BoolCircuit.gatesAcyclic_embedAll}
\leanok
\uses{BoolCircuit.DAGCircuit, BoolCircuit.GatesAcyclic, BoolCircuit.GatesAcyclic.append, BoolCircuit.GatesAcyclic.nil, BoolCircuit.embedAll, BoolCircuit.gatesAcyclic_embedGates}
\difficulty{4}
Let $m,L\in\mathbb{N}$, let $D_1,\ldots,D_r$ be finite directed acyclic circuits with
$m$ input vertices, and let $\mathrm{srcs}$ be a list of $m$ vertex indices, each less
than $L$. Then the list obtained by embedding these circuits in succession, using the
same source list and advancing the offset by each circuit's number of gates plus one, is
acyclic relative to the initial bound $L$.
\end{theorem}

\begin{theorem}[Fan-in after embedding circuits]
\lean{BoolCircuit.faninTwo_embedAll}
\label{BoolCircuit.faninTwo_embedAll}
\leanok
\uses{BoolCircuit.DAGCircuit, BoolCircuit.DAGCircuit.IsFaninTwo, BoolCircuit.embedAll, BoolCircuit.faninTwo_embedGates}
\difficulty{4}
Let $m$ be a natural number, let $D_1,\ldots,D_r$ be a finite list of circuits with $m$
inputs, and suppose each circuit is well formed and every gate in it has at most two
inputs. For any list of source vertices and any natural number $L$, every gate in the
circuit obtained by embedding this list is well formed and has at most two inputs.
\end{theorem}
